\section{Large series prior and theoretical guarantees} \label{sec:method}

Section~\ref{sec:prior-def} defines the prior. Section~\ref{sec:rates-assumptions} introduces the sequences used by the main theorem and states three assumptions, each with a discussion of its role. Section~\ref{sec:contraction} presents the posterior contraction theorem and records, from \citet{WeiGhosal2020}, conditions under which the Horseshoe with a deterministic global scale satisfies the prior conditions. Section~\ref{sec:adaptive-rates} specializes the theorem to H\"older-type classes with the Fourier, Chebyshev, and multiresolution B-spline dictionaries, which yields adaptive contraction rates from one fixed prior; these rates are near-minimax for the classes of the Fourier and B-spline dictionaries and of the univariate Chebyshev dictionary. All proofs and the approximation lemmas that support Section~\ref{sec:adaptive-rates} are in the Appendix.

\subsection{The prior}\label{sec:prior-def}

Throughout, a \emph{dictionary} is a finite collection $\{\xi_1,\ldots,\xi_{J_n}\}$ of known functions used to expand the regression function. It is built from a standard system, such as B-splines, Fourier functions, or the Chebyshev system on $[0,1]$, consisting of the constant $1$ and the functions $\sqrt2\,T_k(2x-1)$, $k\geq1$; its size is deliberately larger than the number of terms needed to approximate $f_0$ well. The enlarged size ensures that an accurate approximation of $f_0$ exists within the dictionary for a wide range of smoothness levels simultaneously; regularization is then delegated to the prior, which suppresses the terms not required by the particular regression function. We use a deterministic but large dimension $J_n\to\infty$; in the principal examples $J_n=\lceil n^b\rceil$ for a constant $b>0$. All subsequent definitions and asymptotic statements are understood for $n$ sufficiently large. For $\theta=(\theta_1,\ldots,\theta_{J_n})\in\mathbb{R}^{J_n}$, set
\begin{align*}
f_\theta(x)=\sum_{j=1}^{J_n}\theta_j\xi_j(x)=\theta^T\xi(x).
\end{align*}
For every $n$, the coefficient prior is the triangular-array product
\begin{align*}
\Pi_n=\bigotimes_{j=1}^{J_n}\Pi_{n,j},
\qquad d\Pi_{n,j}(t)=\pi_n(t)\,dt,
\end{align*}
where the common marginal density $\pi_n$ may depend on $n$. We use $\Pi_n$ also for the induced prior and posterior on $f=f_\theta$. The only structural choices are the dictionary (its family and growth exponent $b$) and the marginal density $\pi_n$: there is no prior on the number of terms, and hence no varying-dimensional computation. The dictionary need not be linearly independent: no result uses identifiability of $\theta$. The contraction statements concern $f_\theta$, and the support statements count coordinates without recovering them. This freedom is what allows the redundant multiresolution B-spline dictionaries of Section~\ref{sec:adaptive-rates}.

\subsection{Rate sequences and assumptions}\label{sec:rates-assumptions}

The results are expressed through four given positive sequences, from which the contraction rate $\epsilon_n$ is derived below:
\begin{itemize}

\item $\en>0$ is the \emph{approximation accuracy}: the accuracy at which $f_0$ is approximated by a sparse element of the dictionary. It is also the scale on which the prior-mass conditions below are expressed.
\item $a_n\in\{1,\ldots,J_n\}$ is the \emph{effective sparsity}: the number of dictionary coefficients allowed to be non-negligible in that approximation.
\item $\delta_n>0$ is the \emph{negligibility threshold}: a coefficient below $\delta_n$ in absolute value is treated as zero. It is taken sufficiently small that $J_n$ coordinates of magnitude $\delta_n$ jointly perturb $f_\theta$ by at most a multiple of $\en$.
\item $E_n>0$ is the \emph{coefficient bound}: the largest magnitude allowed for a coefficient of the approximant. No monotonicity or lower bound is assumed for $E_n$.

\end{itemize}
From these, define
\begin{align}
s_n&=\max\left\{a_n,\left\lceil\frac{n\en^2}{\log J_n}\right\rceil\right\},
&\epsilon_n&=\sqrt{\frac{s_n\log J_n}{n}},
&n\epsilon_n^2&=s_n\log J_n. \label{sparsity}
\end{align}
The quantity $\epsilon_n$ takes the form of the sparse-regression rate: identifying $s_n$ active coefficients among $J_n$ costs $s_n\log J_n/n$ in squared empirical error, to be compared with the squared approximation error $\en^2$. The standing sequence requirements are
\begin{align}
n\en^2\longrightarrow\infty,
\qquad
\delta_n\leq\frac{\en}{J_n},
\qquad
s_n\log J_n=o(n). \label{eq:standing-sequences}
\end{align}
The first requirement ensures that the exponential bounds below, stated on the scale $e^{-n\en^2}$, tend to zero; the second formalizes the smallness of the negligibility threshold described above; the third requires the selection cost $s_n\log J_n$ to be of smaller order than $n$, which yields $\epsilon_n\to0$. Consequently, $\epsilon_n\geq\en$, $n\epsilon_n^2\to\infty$, and $\epsilon_n\to0$, while
\begin{align}
\epsilon_n^2\asymp\max\left\{\en^2,\frac{a_n\log J_n}{n}\right\}. \label{eq:two-rates}
\end{align}
The attained rate is therefore determined by the larger of the approximation error and the sparse-estimation error. In particular, if $n\en^2\lesssim a_n\log J_n$, then $s_n\asymp a_n$ and $\epsilon_n\asymp\sqrt{a_n\log J_n/n}$, so the rate is driven by the effective sparsity $a_n$. This is the regime of all the examples below.

We can now state the three assumptions: one on the dictionary, one on the truth, and one on the prior.

\begin{assumption}[Basis]\label{ass:basis}
There is a constant $K > 0$ such that, for every $n$,
\begin{align}
\max_{1\leq j\leq J_n}\|\xi_j\|_{2,n}\leq K. \label{basis}
\end{align}
\end{assumption}

Assumption~\ref{ass:basis} is a normalization requirement: every dictionary element has empirical norm bounded by a fixed constant. Fourier atoms, B-splines, and Chebyshev atoms on $[0,1]$ satisfy it with explicit constants; for the Chebyshev atoms $1$ and $\sqrt2\,T_k(2x-1)$, $k\geq1$, $K=\sqrt2$ for every configuration of design points in $[0,1]$, because $|T_k|\leq1$ on $[-1,1]$ (Lemma~\ref{lem:chebyshev-profile}), and for their tensor products on $[0,1]^d$, $K=2^{d/2}$ for every configuration of design points in $[0,1]^d$ (Lemma~\ref{lem:chebyshev-aniso-profile}).

\begin{assumption}[Sparse approximability]\label{ass:sparse}
For all sufficiently large $n$, there exists $\theta_0=\theta_0(n)\in\mathbb{R}^{J_n}$ such that
\begin{align}
\|f_0-f_{\theta_0}\|_{2,n}&\leq\en, \label{approx}\\
\#\{j:|\theta_{0,j}|>\delta_n\}&\leq a_n,
&\max_{1\leq j\leq J_n}|\theta_{0,j}|&\leq E_n. \label{eq:sparse-approximant}
\end{align}
For one such approximant, fix a set $\xi_n^*\subset\{1,\ldots,J_n\}$ of exactly $a_n$ indices containing every $j$ for which $|\theta_{0,j}|>\delta_n$. Thus $|\theta_{0,j}|\leq\delta_n$ for every $j\notin\xi_n^*$.
\end{assumption}

Assumption~\ref{ass:sparse} requires an approximant within $\en$ of $f_0$ whose coefficients are bounded by $E_n$ and of which at most $a_n$ exceed the negligibility threshold $\delta_n$. By Assumption~\ref{ass:basis} and~\eqref{eq:standing-sequences}, setting the remaining coefficients to zero changes $f_{\theta_0}$ by at most $KJ_n\delta_n\leq K\en$ in the empirical norm. The index set $\xi_n^*$ records the active coordinates; it is a device of the analysis, not an input to the method. Corollary~\ref{cor:anisotropic} verifies this assumption for H\"older-type $f_0$ with the Fourier and multiresolution B-spline dictionaries.

\begin{assumption}[Prior conditions]\label{ass:prior}
The coordinates are independent and identically distributed under $\Pi_n$ with common density $\pi_n$. Assumption~\ref{ass:sparse} is understood to fix one admissible pair $(\theta_0,\xi_n^*)$; because the coordinate prior is iid and $|\xi_n^*|=a_n$, the probability in~\eqref{eq:4} depends only on $J_n-a_n$, not on the particular admissible set or on the values of $\theta_0$. Fix a constant $\eta\in(0,1)$. The following conditions hold for all sufficiently large $n$:
\begin{align}
\int_{|t|>\delta_n}\pi_n(t)\,dt
&\leq J_n^{-(1+\mu)}
\quad\text{for a constant }\mu>0, \tag{WG3}\label{eq:3}\\
-\log\Pi_n\!\left(\sum_{j\notin\xi_n^*}|\theta_j|\leq\eta\en\right)
&\leq c_4n\en^2
\quad\text{for a constant }c_4\geq0. \tag{WG4}\label{eq:4}
\end{align}
Let
\begin{align*}
\pinf=\inf_{|t|\leq E_n}\pi_n(t)>0.
\end{align*}
For a constant $c_5\geq0$, assume
\begin{align}
a_n\left[\log a_n+2\log\left(\frac1\en\right)-\log\pinf\right]
\leq c_5n\en^2. \tag{WG5}\label{eq:5}
\end{align}
\end{assumption}

The three conditions play distinct roles. Condition~\eqref{eq:3} requires each individual coordinate to be negligible with probability at least $1-J_n^{-(1+\mu)}$. The prior expected number of coordinates exceeding $\delta_n$ is then at most $J_n^{-\mu}\to0$, so a priori almost all mass lies near the zero function, in accordance with the sparsity of the approximant in Assumption~\ref{ass:sparse}. Condition~\eqref{eq:4} requires that the event that all $J_n-a_n$ inactive coordinates are jointly small has probability at least $e^{-c_4n\en^2}$. Condition~\eqref{eq:5} requires the marginal density to retain enough mass on $[-E_n,E_n]$ that placing the $a_n$ active coordinates near their required values costs prior probability only $e^{-O(n\en^2)}$, a cost the likelihood overcomes. In short, the prior must be concentrated enough to suppress the $J_n-a_n$ terms that $f_0$ does not need, and spread out enough to reach the $a_n$ terms that it does. The three conditions adapt conditions (3)--(5) of \citet{WeiGhosal2020}, stated there for high-dimensional logistic regression, with their dimension $p$ replaced by $J_n$; the labels~\eqref{eq:3}--\eqref{eq:5} recall that numbering.

\subsection{Posterior contraction}\label{sec:contraction}

The main theorem states that, under the above assumptions, the posterior concentrates around $f_0$ at the rate $\epsilon_n$ and, at the same time, concentrates on coefficient vectors with at most a bounded multiple of $s_n$ coordinates exceeding $\delta_n$ in absolute value.

\begin{theorem}\label{thm:main}
Suppose model~\eqref{eq:model}, Assumptions~\ref{ass:basis}--\ref{ass:prior}, and the standing requirements~\eqref{eq:standing-sequences} hold, and let $M_n\to\infty$ be an arbitrary sequence. Then there exists a fixed integer $A\geq1$, depending only on $(K,\sigma,\mu,\eta,c_4,c_5)$, such that, for all sufficiently large $n$,
\begin{align}
E_0\,\Pi_n\!\left(\|f_0-f\|_{2,n}>M_n\epsilon_n\mid Y^{(n)}\right)
\leq6e^{-n\en^2}. \label{eq:main-L1}
\end{align}
Consequently,
\begin{align}
\Pi_n\!\left(\|f_0-f\|_{2,n}>M_n\epsilon_n\mid Y^{(n)}\right)
\xrightarrow{P_0}0. \label{eq:main-contraction}
\end{align}
Moreover,
\begin{align}
E_0\,\Pi_n\!\left(\#\{j:|\theta_j|>\delta_n\}>As_n\mid Y^{(n)}\right)
\leq2e^{-n\en^2}, \label{eq:main-support-L1}
\end{align}
and the corresponding posterior probability converges to zero in $P_0$-probability.
\end{theorem}

Theorem~\ref{thm:main} makes two assertions. The first, \eqref{eq:main-L1}--\eqref{eq:main-contraction}, concerns recovery of the function: apart from posterior mass that is exponentially small in expectation, posterior draws lie within $M_n\epsilon_n$ of $f_0$ in the empirical prediction norm. The second, \eqref{eq:main-support-L1}, provides automatic complexity control: the posterior concentrates on coefficient vectors with at most $As_n$ non-negligible coordinates. When $s_n=o(J_n)$, the posterior thus uses a vanishing fraction of the $J_n$ coordinates; this complexity control replaces a hyperprior on the number of terms, without varying-dimensional sampling. Because it counts coordinates rather than identifying them, it holds for redundant dictionaries as well.

The proof, deferred to the Appendix, follows the testing approach of \citet{Ghosal2000} and \citet{Ghosal2007} and rests on four estimates, stated there as Lemmas~\ref{lem:local-mass}--\ref{lem:tests}: a lower bound on the prior mass of a prediction ball around $f_0$, driven by~\eqref{eq:4}--\eqref{eq:5}; a lower bound on the marginal likelihood, obtained from the prior-mass bound through an exact Kullback--Leibler computation for the Gaussian model; a bound showing that the prior places exponentially little mass outside the sparse coefficient vectors, driven by~\eqref{eq:3}; and a test that separates $f_0$ from the sparse vectors at prediction distance beyond $M_n\epsilon_n$, whose construction exploits the fact that such vectors lie, up to negligible coordinates, in low-dimensional evaluation subspaces that admit coverings of controlled cardinality. The theorem follows from the four estimates by the standard decomposition of the posterior probability. Throughout, $\en$ governs the local prior mass and the marginal likelihood, whereas $\epsilon_n$ governs the tests and sets the radius of the contraction statement.

Similar theoretical guarantees exist in related settings. \citet{WeiGhosal2020} established analogous contraction properties for shrinkage priors in high-dimensional logistic regression, while \citet{SongLiang2023} analyzed high-dimensional linear regression with a number of covariates that may substantially exceed the sample size. Our setting is instead nonparametric: the sparse vector is an approximating series representation of $f_0$, and Assumption~\ref{ass:sparse} makes this structural requirement explicit.

A concrete prior satisfying Assumption~\ref{ass:prior} is the Horseshoe \citep{Carvalho2009,Carvalho2010}, $\theta_j\mid\lambda_j\sim N(0,\tau^2\lambda_j^2)$, $\lambda_j\overset{\mathrm{iid}}{\sim}C^+(0,1)$, with a deterministic global scale $\tau=\tau_n$, as in \citet{VanDerPas2014}; the coordinates are then iid. Map the notation $(p,s_*,a_n,\bar\epsilon_n,E,\alpha)$ of \citet{WeiGhosal2020} to $(J_n,a_n,\delta_n,\en,E_n,\tau_n)$. If, for a constant $\mu'>0$,
\begin{gather}
\tau_n\leq\frac{J_n^{-(2+\mu')}}{\sqrt n},\qquad
-\log\tau_n=O(\log J_n),\qquad
\log E_n=O(\log J_n),\notag\\
\delta_n\asymp\frac{\en}{J_n},\qquad
a_n\log J_n\lesssim n\en^2,\qquad
\frac{n\en^2}{\log J_n}\to\infty, \label{eq:hs-conditions}
\end{gather}
then the bounds in the proof of their Theorem~3.3 give Assumption~\ref{ass:prior}, with any $\mu<\mu'$ in~\eqref{eq:3}. Their condition~(6), a tail condition on the prior used to truncate large coefficients, is not needed, because the tests in the proof of Theorem~\ref{thm:main} depend on $\theta$ only through $\|f_\theta-f_0\|_{2,n}$. For $\tau_n=J_n^{-(2+\mu')}/\sqrt n$ and the sequences of Corollary~\ref{cor:anisotropic}, all conditions in~\eqref{eq:hs-conditions} hold for every smoothness vector: $J_n$ grows polynomially in $n$, so $-\log\tau_n=O(\log J_n)$; $E_n=E$ is constant, $\delta_n=c_\delta\en/J_n$, $a_n\log J_n\asymp n\en^2$, and $n\en^2/\log J_n\asymp a_n\to\infty$. This scale depends only on $n$ and $J_n$, not on the smoothness of $f_0$, which is what makes the adaptation of Section~\ref{sec:adaptive-rates} possible.

\subsection{Adaptive contraction rates}\label{sec:adaptive-rates}

Theorem~\ref{thm:main} converts sparse approximability of $f_0$ into a contraction rate. To obtain rates over classical smoothness classes, we therefore need one dictionary-specific input, which we call a \emph{sparse approximation profile}: a family of coefficient vectors, indexed by an integer term budget, whose supports grow proportionally to the budget while their approximation errors decay polynomially in it. Specifically, a budget of $m_k$ terms in coordinate direction $k$ yields accuracy of order $m_k^{-\alpha_k}$, with support of order $\prod_k m_k$ and uniformly bounded coefficients.

The three dictionaries treated in this paper, the multiresolution B-spline dictionary, the Fourier dictionary, and the Chebyshev dictionary, have such a profile; see Lemmas~\ref{lem:bspline-profile}, \ref{lem:fourier-profile}, \ref{lem:chebyshev-profile}, and~\ref{lem:chebyshev-aniso-profile} in the Appendix. For B-splines, the profile builds on the anisotropic approximation result of \citet[Lemma~2]{ShenGhosal2015}. For the Chebyshev dictionary and its tensor products, the profile is established in every dimension, at every directional smoothness $\alpha_k>0$ and with no periodicity or boundary condition, for functions whose univariate sections along each coordinate are H\"older of order $\alpha_k$ uniformly over the other coordinates: the change of variables $x_k=\{1+\cos(2\pi u_k)\}/2$ in each coordinate turns the Chebyshev atoms into products of cosines and reduces the proof to the Fourier case.

In the univariate case, the corollary below states that if $f_0$ lies in a H\"older ball $C^\alpha([0,1])$, the dictionary size satisfies $J_n\asymp n^b$ with $b\geq1/(2\alpha+1)$, and the sparse approximation profile and prior conditions stated there hold, then the posterior contracts at the rate $\epsilon_n\asymp(n/\log n)^{-\alpha/(2\alpha+1)}$, the minimax rate up to a logarithmic factor, even though neither the prior nor the dictionary depends on $\alpha$. The rate is reached through sparsity: the posterior concentrates on coefficient vectors with at most a fixed multiple of $(n/\log n)^{1/(2\alpha+1)}$ non-negligible coordinates. The general statement permits distinct smoothness levels $\alpha_1,\ldots,\alpha_d$ in different coordinate directions, in which case the rate is governed by the harmonic mean $\alpha^*$ of the directional smoothness parameters.

\begin{corollary}{\thecorollary}\label{cor:anisotropic}
Let $d$ be fixed, let $R>0$, and let
$\boldsymbol\alpha=(\alpha_1,\ldots,\alpha_d)\in(0,\infty)^d$.
Assume that $f_0$ lies in a ball of radius $R$ of an anisotropic H\"older-type class $C^{\boldsymbol\alpha}([0,1]^d)$; the precise classes for the three dictionaries of this paper are specified in Lemmas~\ref{lem:bspline-profile}, \ref{lem:fourier-profile}, \ref{lem:chebyshev-profile}, and~\ref{lem:chebyshev-aniso-profile}; for the tensor-product Chebyshev dictionary with $d>1$ the class is the sectional class $C^{\boldsymbol\alpha}_{\mathrm{sec}}([0,1]^d;R)$ of Lemma~\ref{lem:chebyshev-aniso-profile}, the functions whose univariate sections in each direction $k$ lie in the ball of radius $R$ of $C^{\alpha_k}([0,1])$, and the smoothness enters the result only through the sparse approximation profile below. Define
\begin{align}
\alpha^*
&=\frac{d}{\sum_{k=1}^d\alpha_k^{-1}},
&
r
&=\frac{\alpha^*}{2\alpha^*+d},
&
D
&=\frac{d}{2\alpha^*+d}, \label{eq:aniso-exponents}\\
\beta_k
&=\frac{\alpha^*}{\alpha_k(2\alpha^*+d)},
\qquad k=1,\ldots,d. \label{eq:aniso-beta}
\end{align}
Let the dictionary be a tensor product with coordinate sizes $L_{n,k}\asymp n^{b_k}$ for fixed constants $b_k\geq\beta_k$ and total size $J_n\asymp\prod_{k=1}^dL_{n,k}$. Suppose Assumption~\ref{ass:basis} holds and that $f_0$ admits the following sparse approximation profile: for every integer rectangle $\boldsymbol m=(m_1,\ldots,m_d)$ with $1\leq m_k\leq L_{n,k}$, there is a vector $\theta(\boldsymbol m)\in\R^{J_n}$ satisfying
\begin{align}
\#\operatorname{supp}\theta(\boldsymbol m)
&\leq C_{\mathrm{sp}}\prod_{k=1}^dm_k,\notag\\
\|f_0-f_{\theta(\boldsymbol m)}\|_{2,n}
&\leq C_{\mathrm{app}}\sum_{k=1}^dm_k^{-\alpha_k}, \label{eq:aniso-profile}\\
\|\theta(\boldsymbol m)\|_\infty
&\leq E.\notag
\end{align}
Here $C_{\mathrm{sp}},C_{\mathrm{app}},E$ are fixed positive constants independent of $n$ and $\boldsymbol m$, although they may depend on $(d,R,\boldsymbol\alpha)$ and on fixed basis orders. Lemmas~\ref{lem:bspline-profile}, \ref{lem:fourier-profile}, \ref{lem:chebyshev-profile}, and~\ref{lem:chebyshev-aniso-profile} supply this profile for the multiresolution B-spline dictionary, under the spline-order restriction and, in the anisotropic case, the integer-smoothness qualification stated there; for the Fourier dictionary when $f_0$ is periodic; and for the tensor-product Chebyshev dictionary, in every dimension $d\geq1$ and for every $\boldsymbol\alpha\in(0,\infty)^d$, when the design points lie in $[0,1]^d$ and $f_0$ lies in the sectional class $C^{\boldsymbol\alpha}_{\mathrm{sec}}([0,1]^d;R)$, with $C_{\mathrm{sp}}=1$ and $E=R$. The sectional class contains, with the same radius, the isotropic H\"older ball on $[0,1]^d$ and the integer anisotropic class used for splines, and for $d=1$ it is the ball of radius $R$ of $C^\alpha([0,1])$. Given the profile, fix constants $C_\epsilon\geq d\,C_{\mathrm{app}}$ and $c_\delta\in(0,1)$, and choose the term budget $\boldsymbol m_n=(m_{n,1},\ldots,m_{n,d})$ and the analytical sequences of Section~\ref{sec:rates-assumptions} as
\begin{align}
N_n&=\frac{n}{\log J_n},
&
m_{n,k}&=\left\lceil N_n^{\beta_k}\right\rceil,
&
a_n&=\left\lceil C_{\mathrm{sp}}\prod_{k=1}^dm_{n,k}\right\rceil, \label{eq:aniso-choice}\\
\en&=C_\epsilon N_n^{-r},
&
E_n&=E,
&
\delta_n&=c_\delta\frac{\en}{J_n}. \label{eq:aniso-prior-choice}
\end{align}
The budget balances the squared approximation error against the sparse-estimation penalty, as explained in the Appendix. Suppose the coefficient prior satisfies Assumption~\ref{ass:prior} for these sequences, and let $s_n$ and $\epsilon_n$ be defined by~\eqref{sparsity}. Then Assumption~\ref{ass:sparse} and the standing requirements~\eqref{eq:standing-sequences} hold, so that Theorem~\ref{thm:main} applies, and:
\begin{enumerate}[label=(\roman*)]
\item $s_n\asymp a_n\asymp(n/\log n)^{d/(2\alpha^*+d)}$ and $\epsilon_n\asymp(n/\log n)^{-\alpha^*/(2\alpha^*+d)}$;
\item\label{cor:anisotropic-tails} there are constants $A',c>0$, depending only on the integer $A$ of Theorem~\ref{thm:main}, $C_{\mathrm{sp}}$, $C_\epsilon$, and $d$, such that, for every sequence $M_n\to\infty$ and all sufficiently large $n$,
\begin{align}
E_0\,\Pi_n\left(
\|f-f_0\|_{2,n}>M_n\epsilon_n
\,\middle|\,Y^{(n)}
\right)
&\leq6e^{-c n\epsilon_n^2}, \label{eq:cor1-L1}\\
E_0\,\Pi_n\left(
\#\{j:|\theta_j|>\delta_n\}>A'a_n
\,\middle|\,Y^{(n)}
\right)
&\leq2e^{-c n\epsilon_n^2}; \label{eq:cor1-support}
\end{align}
\item both posterior probabilities in~\ref{cor:anisotropic-tails} converge to zero in $P_0$-probability.
\end{enumerate}
\end{corollary}

Part~\ref{cor:anisotropic-tails} retains the contraction radius $M_n\epsilon_n$ of Theorem~\ref{thm:main}. For the balanced sequences of this corollary, $s_n\asymp a_n$ and $\en\asymp\epsilon_n$, so the support bound can be stated in terms of $a_n$ and the exponential bounds in terms of $n\epsilon_n^2$, with fixed constants $A'$ and $c$.

The sequences~\eqref{eq:aniso-choice}--\eqref{eq:aniso-prior-choice} are analytical devices of the proof and depend on $\boldsymbol\alpha$; the prior and the dictionary do not. Adaptation over a set of smoothness vectors therefore means that one fixed dictionary satisfies $b_k\geq\beta_k$ for every $\boldsymbol\alpha$ in the set and one fixed prior satisfies Assumption~\ref{ass:prior} for the sequences of every $\boldsymbol\alpha$ in the set. The adaptation is pointwise: for each fixed truth and smoothness vector, the rate is attained in the empirical prediction norm. The resulting admissible ranges are described in the Appendix. The Horseshoe of Section~\ref{sec:contraction} meets the prior requirement: \eqref{eq:hs-conditions} holds for the sequences of every $\boldsymbol\alpha$, and its global scale depends only on $n$ and $J_n$. One Horseshoe prior, fixed before the data are seen, therefore attains the stated contraction rate pointwise throughout the admissible range; for the classes of Lemmas~\ref{lem:bspline-profile}, \ref{lem:fourier-profile}, and~\ref{lem:chebyshev-profile} this is the near-minimax rate.

Corollary~\ref{cor:anisotropic} contains the isotropic and univariate settings verbatim as special cases; no separate statements or proofs are needed. If $\alpha_1=\cdots=\alpha_d=\alpha$, then $\alpha^*=\alpha$, $r=\alpha/(2\alpha+d)$, and $\beta_k=1/(2\alpha+d)$ for every $k$; all coordinate sizes may be taken equal, $L_{n,k}=L_n\asymp n^b$ with $b\geq1/(2\alpha+d)$, so that the total size is $J_n\asymp n^{db}$. The budget is then the cube $m_{n,k}=\lceil N_n^{1/(2\alpha+d)}\rceil$, and the corollary gives $a_n\asymp(n/\log n)^{d/(2\alpha+d)}$ and $\epsilon_n\asymp(n/\log n)^{-\alpha/(2\alpha+d)}$. If moreover $d=1$, then $J_n=L_n\asymp n^b$ with $b\geq1/(2\alpha+1)$, $a_n\asymp(n/\log n)^{1/(2\alpha+1)}$, and $\epsilon_n\asymp(n/\log n)^{-\alpha/(2\alpha+1)}$.

\begin{remark}[General polynomial choices]\label{rem:polynomial}
The choices~\eqref{eq:aniso-choice}--\eqref{eq:aniso-prior-choice} balance approximation error and selection cost in Theorem~\ref{thm:main}. More generally, under the hypotheses of Theorem~\ref{thm:main}, if $J_n\asymp n^b$, $a_n\asymp n^\gamma$ with $\gamma<1$, and $n\en^2\lesssim a_n\log J_n$, then $s_n\asymp a_n$ and $\epsilon_n\asymp\sqrt{\log n}\,n^{(\gamma-1)/2}$. The exponent $\gamma$ measures how sparsely the basis system represents $f_0$, and it is the dominant factor in the rate: the sparser the representation, the faster the contraction. Since $\gamma$ may differ greatly from one basis system to another for the same $f_0$, a dictionary matters mainly through the sparsity it induces; Section~\ref{sec:convergence} shows this effect empirically.
\end{remark}

\begin{remark}[Practical guidance]\label{rem:practical}
The dictionary matters through its sparse approximation profile, as described in the previous remark. A single-resolution B-spline basis has no sparse profile for functions bounded away from zero (see the Appendix), which is why the B-spline case of Corollary~\ref{cor:anisotropic} uses a multiresolution dictionary; global systems such as the Fourier dictionary on its periodic class and the tensor-product Chebyshev dictionary on its sectional class have one without any multiresolution construction. As for the growth exponent, $b=1$ covers every fixed smoothness level as far as dictionary growth is concerned. In the univariate case, $b=0.7$ and $b=0.8$ give smaller dictionaries and satisfy the growth requirement for $\alpha\geq3/14$ and $\alpha\geq1/8$, respectively; the dictionary-specific smoothness restrictions still apply. The Horseshoe result of Section~\ref{sec:contraction} assumes a deterministic global scale $\tau_n$; in practice, as is standard since the original proposal of \citet{Carvalho2010}, the global scale is given a half-Cauchy prior, as in Sections~\ref{sec:recovery}--\ref{sec:convergence}.
\end{remark}
